\subsection{Finite bases and cost repairs}
\label{sec:special-cases}

These are not extra algebraic existence bases.  Degree~$15$ is the optimized,
fused $k=1$ endpoint of the $8k+7$ branch: it creates the quartic missing
from the one-product degree-$3$ pair without exceeding seven products.
Degrees~$27$ and~$31$ are the two bridges around the unavailable
three-product degree-$7$ joint realization.  Their common top-down decoder is
displayed after the circuits.

\subsubsection{Degree $3$}
\begin{lemma}[A one-product compatible realization for $3$]
    \label{lem:base-three-compatible}
    There exists a splittable pair for $3$ that is compatible on $\rng 3$,
    has a one-realization, and records its monic quadratic $H_2$.
\end{lemma}
\begin{proof}
    Define $H_2[\alpha_1,\alpha_2](x)=(x+\alpha_2)x+\alpha_1$ and set
    \[
      T^{(1)}_3[\alpha_0,\alpha_1,\alpha_2](x)\coloneqq H_2,\qquad
      T^{(2)}_3[\alpha_0,\alpha_1,\alpha_2](x)\coloneqq H_2+\alpha_0.
    \]
    Then $T^{(1)}_3$ and $T^{(2)}_3$ are monic of degree $2$ and
    \[
      P_3=xT^{(1)}_3+T^{(2)}_3
        =x^3+(\alpha_2+1)x^2+(\alpha_1+\alpha_2)x+(\alpha_0+\alpha_1).
    \]
    Hence $\alpha_2=[x^2]P_3-1$, $\alpha_1=[x^1]P_3-\alpha_2$, and
    $\alpha_0=[x^0]P_3-\alpha_1$, so $P_3$ is decodable.  We also verify
    the causal window needed by the induction.  For
    $\Phi_3=xT^{(1)}_3+T^{(2)}_3=P_3$, the relevant coefficients are

    \[
      \begin{array}{lll}
        [x^2]T^{(1)}_3=[x^2]T^{(2)}_3=1,&
        [x^1]T^{(1)}_3=[x^1]T^{(2)}_3=\Phi_{3,2}-1,\\[2pt]
        [x^0]T^{(1)}_3=\Phi_{3,1}-\Phi_{3,2}+1,&
        [x^0]T^{(2)}_3=\Phi_{3,0},
      \end{array}
    \]
    where $\Phi_{3,j}=[x^j]\Phi_3$.  The first-component formulas use
    only rows $\ge j+1$, and the second-component formulas only rows
    $\ge j$, exactly as required in \Cref{def:compatible-pair}.  Thus
    $(T^{(1)}_3,T^{(2)}_3)$ is compatible on $\rng 3$, and in particular it
    is splittable.  The single displayed product computes $H_2$; both
    components are then obtained by additions, so this is a one-realization
    and $H_2$ is a recorded byproduct.
\end{proof}

\subsubsection{Degree $7$: direct septic base}
The canonical construction in \Cref{rem:canonical-split} makes the septic
algebraically splittable.  What the optimal recursion would need, but we do
not have, is a compatible joint realization of its two components using
three multiplications and recording the required power byproducts.  A direct
four-multiplication evaluation scheme is enough for the final polynomial,
and the recursion below is arranged never to call degree~$7$.

\begin{lemma}[A decodable septic in characteristic different from $2$]
\label{lem:septic-base}
Assume $\operatorname{char}(\mathbb F)\ne2$.  Define
\begin{align}
  y&=x(x+\alpha_6),\\
  z&=(\alpha_5+x+y)(\alpha_4+x),\\
  w&=(\alpha_3+z)x,\\
  v&=(\alpha_2+x+z)(\alpha_1+w),\\
  P_7&=\alpha_0+y+w+v.
\end{align}
Then $P_7$ is monic of degree~$7$, uses four multiplications, and is
decodable.
\end{lemma}

\begin{figure}[H]
  \centering
  \resizebox{0.95\textwidth}{!}{% Circuit for the degree-7 base construction (lem:septic-base): four multiplications.
%   y = x(x+a6),  z = (a5+x+y)(a4+x),  w = (a3+z)x,  v = (a2+x+z)(a1+w),  P7 = a0+y+w+v
\begin{tikzpicture}[
    >=Latex, thick,
    mul/.style={circle, draw, inner sep=1.5pt, minimum size=6.5mm, font=\large, fill=white},
    inp/.style={font=\small},
    lbl/.style={font=\footnotesize, inner sep=1pt, fill=white, fill opacity=0.85, text opacity=1}
  ]
  \node[inp] (x) at (0,0) {$x$};
  \node[mul] (my) at (2.2,0) {$\times$};
  \node[mul] (mz) at (4.9,0) {$\times$};
  \node[mul] (mw) at (7.6,0) {$\times$};
  \node[mul] (mv) at (10.3,0) {$\times$};
  \node[inp] (out) at (13.6,0) {$P_7$};
  \node[font=\footnotesize] at (2.2,0.55) {$y$};
  \node[font=\footnotesize] at (4.9,0.55) {$z$};
  \node[font=\footnotesize] at (7.6,0.55) {$w$};
  \node[font=\footnotesize] at (10.3,0.55) {$v$};

  % y = x (x + a6)
  \draw[->] (x) to[bend left=30] node[lbl, above] {$x$} (my);
  \draw[->] (x) to[bend right=30] node[lbl, below] {$x+\alpha_6$} (my);
  % z = (a5 + x + y)(a4 + x)
  \draw[->] (my) to[bend left=30] node[lbl, above] {$\alpha_5+x+y$} (mz);
  \draw[->] (x) to[out=-60, in=-115, looseness=1.0] node[lbl, pos=0.78, below] {$\alpha_4+x$} (mz);
  % w = (a3 + z) x
  \draw[->] (mz) to[bend left=30] node[lbl, above] {$\alpha_3+z$} (mw);
  \draw[->] (x) to[out=-80, in=-105, looseness=1.05] node[lbl, pos=0.5, below] {$x$} (mw);
  % v = (a2 + x + z)(a1 + w)
  \draw[->] (mw) to[bend left=30] node[lbl, above] {$\alpha_1+w$} (mv);
  \draw[->] (mz) to[out=-45, in=-135, looseness=1.0] node[lbl, pos=0.55, below] {$\alpha_2+x+z$} (mv);
  % output: a0 + y + w + v
  \draw[->] (mv) -- node[lbl, above] {$\alpha_0+y+w+v$} (out);
  % Use the visible ink extent, not the lower Bezier control points.
  % This changes layout only: no clipping path is installed.
  \pgfresetboundingbox
  \path[use as bounding box] (-0.2,-2.9) rectangle (13.9,0.8);
\end{tikzpicture}
}
  \caption{The four-multiplication circuit for the septic base $P_7$ of \Cref{lem:septic-base}.
  Every multiplication gate takes two affine combinations of $x$ and previously computed values; the seven parameters $\alpha_0,\ldots,\alpha_6$ enter only as constant terms of these affine forms.}
  \label{fig:septic-circuit}
\end{figure}

\begin{proof}
\Cref{fig:septic-circuit} draws the circuit.
Write
\[
  P_7=x^7+c_6x^6+c_5x^5+c_4x^4+c_3x^3+c_2x^2+c_1x+c_0
\]
and write $z=x^3+z_2x^2+z_1x+z_0$.  The following triangular procedure
recovers the parameters from $(c_0,\ldots,c_6)$:
\begin{align}
 z_2&=c_6/2,\\
 z_1&=(c_5-z_2^2-1)/2,\\
 R&=c_4-1-2z_2z_1-z_2,\\
 \alpha_1&=c_3-z_2-z_2R-z_1^2-z_1,\\
 W&=c_2-z_1-1-z_2\alpha_1-z_1R,\\
 \alpha_6&=c_1-(z_1+1)\alpha_1-W(R+1-W),\\
 \alpha_4&=z_2-1-\alpha_6,\\
 \alpha_5&=z_1-\alpha_4(1+\alpha_6),\\
 z_0&=\alpha_4\alpha_5,\\
 \alpha_3&=W-z_0,\\
 \alpha_2&=R-2z_0-\alpha_3,\\
 \alpha_0&=c_0-(z_0+\alpha_2)\alpha_1.
\end{align}
Indeed, direct expansion gives
\[
 z_2=\alpha_4+1+\alpha_6,\qquad
 z_1=\alpha_4(1+\alpha_6)+\alpha_5,\qquad
 z_0=\alpha_4\alpha_5,
\]
while the coefficient equations for $P_7$ give, successively,
$R=2z_0+\alpha_2+\alpha_3$ and $W=z_0+\alpha_3$ and then the displayed
formulas for $\alpha_1$ and $\alpha_6$.
Substitution recovers every parameter.  The only division is by~$2$, which
is valid under the characteristic assumption.
\end{proof}


\subsubsection{Degree $15$}
\begin{algorithm}[H]
    \caption{
        Seven-product compatible joint realization in degree $15$.
    }

    \begin{algorithmic}
        \State $H_2[\alpha_6, \alpha_7](x) = (x + \alpha_7)x + \alpha_6$
        \State $H_4[\alpha_4, \alpha_5, \alpha_6, \alpha_7](x) = (H_2 + (x + \alpha_5))(H_2 - (x + \alpha_5)) + \alpha_4$
        \\
        \State $S^{(1)}_1 = Q_7[\alpha_8, \ldots, \alpha_{14}](x, H_2, H_4)$
        \State $S^{(1)}_2 = H_2 + \alpha_{3}$
        \State $S^{(1)}_3 = \alpha_1$
        \State $T^{(1)}_{15}[\alpha_0, \ldots, \alpha_{14}](x)
            = (S^{(1)}_1 + S^{(1)}_2)(S^{(1)}_1 - S^{(1)}_2) + S^{(1)}_3
            = (S^{(1)}_1)^2 - (S^{(1)}_2)^2 + S^{(1)}_3$
        \\
        \State $S^{(2)}_1 = H_4$
        \State $S^{(2)}_2 = H_2 + \alpha_2$
        \State $S^{(2)}_3 = \alpha_0$
        \State $T^{(2)}_{15}[\alpha_0, \ldots, \alpha_{14}](x)
            = (S^{(2)}_1 + S^{(2)}_2)(S^{(2)}_1 - S^{(2)}_2) + S^{(2)}_3 + T^{(1)}_{15}
            = (S^{(2)}_1)^2 - (S^{(2)}_2)^2 + S^{(2)}_3 + T^{(1)}_{15}$

        \State \begin{align}
            P_{15}[\alpha_0, \ldots, \alpha_{14}](x)
            &= x T^{(1)}_{15} + T^{(2)}_{15}
            \\&= (x+1)T^{(1)}_{15} + H_4^2 - (H_2 + \alpha_2)^2 + \alpha_0
        \end{align}
    \end{algorithmic}
\end{algorithm}

\subsubsection{Degree $27$}
\begin{algorithm}[H]
    \caption{
        Thirteen-product compatible joint realization in degree $27$.
    }

    \begin{algorithmic}
        \State $H_2[\alpha_2, \alpha_3](x) = (x + \alpha_3)x + \alpha_2$
        \\
        \State $S^{(1)}_1 = Q_{13}[\alpha_{14}, \ldots, \alpha_{26}](x, H_2)$ \Comment{This call records the quartic $H_4$.}
        \State $S^{(1)}_2 = Q_3[\alpha_4, \ldots, \alpha_6](x, H_2)$
        \State $S^{(1)}_3 = \alpha_1$
        \State $T^{(1)}_{27}[\alpha_0, \ldots, \alpha_{26}](x)
            = (S^{(1)}_1 + S^{(1)}_2)(S^{(1)}_1 - S^{(1)}_2) + S^{(1)}_3
            = (S^{(1)}_1)^2 - (S^{(1)}_2)^2 + S^{(1)}_3$
        \\
        \State $S^{(2)}_1 = Q_7[\alpha_7, \ldots, \alpha_{13}](x, H_2, H_4)$
        \State $S^{(2)}_2 = H_2$
        \State $S^{(2)}_3 = \alpha_0$
        \State $T^{(2)}_{27}[\alpha_0, \ldots, \alpha_{26}](x)
            = (S^{(2)}_1 + S^{(2)}_2)(S^{(2)}_1 - S^{(2)}_2) + S^{(2)}_3 + T^{(1)}_{27}
            = (S^{(2)}_1)^2 - (S^{(2)}_2)^2 + S^{(2)}_3 + T^{(1)}_{27}$

        \State \begin{align}
            P_{27}[\alpha_0, \ldots, \alpha_{26}](x)
            &= x T^{(1)}_{27} + T^{(2)}_{27}
            \\&= (x+1)T^{(1)}_{27} + (S^{(2)}_1)^2 - H_2^2 + \alpha_0
        \end{align}
    \end{algorithmic}
\end{algorithm}

\subsubsection{Degree $31$}
\begin{algorithm}[H]
    \caption{
        Fifteen-product compatible joint realization in degree $31$.
    }

    \begin{algorithmic}
        \State $H_2[\alpha_6, \alpha_7](x) = (x + \alpha_7)x + \alpha_6$
        \State $H_4[\alpha_4, \alpha_5, \alpha_6, \alpha_7](x) = (H_2 + (x + \alpha_5))(H_2 - (x + \alpha_5)) + \alpha_4$
        \\
        \State $S^{(1)}_1 = \bar Q_{15}[\alpha_{16}, \ldots, \alpha_{30}](x, H_2, H_4)$
        \State $S^{(1)}_2 = Q_7[\alpha_8, \ldots, \alpha_{14}](x, H_2, H_4)$
        \State $S^{(1)}_3 = Q_3[\alpha_1, \alpha_2, \alpha_3](x, H_2)$
        \State $T^{(1)}_{31}[\alpha_0, \ldots, \alpha_{30}](x)
            = (S^{(1)}_1 + S^{(1)}_2)(S^{(1)}_1 - S^{(1)}_2) + S^{(1)}_3
            = (S^{(1)}_1)^2 - (S^{(1)}_2)^2 + S^{(1)}_3$
        \\
        \State $S^{(2)}_1 = S^{(1)}_1 + \alpha_{15}$
        \State $S^{(2)}_2 = H_4$
        \State $S^{(2)}_3 = \alpha_0$
        \State $T^{(2)}_{31}[\alpha_0, \ldots, \alpha_{30}](x)
            = (S^{(2)}_1 + S^{(2)}_2)(S^{(2)}_1 - S^{(2)}_2) + S^{(2)}_3
            = (S^{(2)}_1)^2 - (S^{(2)}_2)^2 + S^{(2)}_3$

        \State \begin{align}
            P_{31}[\alpha_0, \ldots, \alpha_{30}](x)
            &= x T^{(1)}_{31} + T^{(2)}_{31}
            \\&= x(S^{(1)}_1)^2 + (S^{(1)}_1 + \alpha_{15})^2 - x(S^{(1)}_2)^2 - H_4^2 + x S^{(1)}_3 + \alpha_0
        \end{align}
    \end{algorithmic}
\end{algorithm}

\begin{figure}[H]
  \centering
  \resizebox{\textwidth}{!}{% Shell decomposition of the three finite splittable constructions.
\begin{tikzpicture}[x=1cm,y=1cm,>=Latex]
  \draw[fp decode] (0,6.25) -- node[fp label, above] {descending coefficient degree}
    (13.55,6.25);
  \node[fp label, anchor=east] at (0,6.25) {high};
  \node[fp label, anchor=west] at (13.55,6.25) {low};

  \node[fp panel title] at (0,5.72) {$n=15$};
  \node[fp active, minimum width=38mm, anchor=west] (a15) at (0,5.02)
    {recover $A=Q_7$\\from $(x+1)A^2$};
  \node[fp seam, minimum width=8mm, anchor=west] at (3.8,5.02) {$+1$};
  \node[fp second, minimum width=39mm, anchor=west] (h15) at (4.6,5.02)
    {recover $H_4$, then $H_2$\\from four rows};
  \node[fp seam, minimum width=17mm, anchor=west] at (8.5,5.02)
    {$-1$\\$-(2b+2)$};
  \node[fp recursive, minimum width=30mm, anchor=west] (l15) at (10.2,5.02)
    {solve $U,V,\alpha_1,\alpha_0$\\from the bottom rows};
  \node[fp label, anchor=west] at (0.35,4.4)
    {decode the parameters inside $A$ only after $H_2,H_4$ are known};

  \node[fp panel title] at (0,3.87) {$n=27$};
  \node[fp active, minimum width=31mm, anchor=west] (a27) at (0,3.17)
    {$A=Q_{13}$\\$(x+1)A^2$};
  \node[fp seam, minimum width=8mm, anchor=west] at (3.1,3.17) {$+1$};
  \node[fp second, minimum width=25mm, anchor=west] (c27) at (3.9,3.17)
    {$C=Q_7$\\$C^2$};
  \node[fp seam, minimum width=8mm, anchor=west] at (6.4,3.17) {$-1$};
  \node[fp pivot, minimum width=25mm, anchor=west] (b27) at (7.2,3.17)
    {$B=Q_3$\\$(x+1)B^2$};
  \node[fp seam, minimum width=8mm, anchor=west] at (9.7,3.17) {$+1$};
  \node[fp recursive, minimum width=30mm, anchor=west] (h27) at (10.5,3.17)
    {recover $H_2,s,r$\\then decode $A,C,B$};
  \draw[fp dependence] (h27.south) to[bend left=16]
    node[fp label, below] {$A$ supplies $H_4$ for $C$} (c27.south);

  \node[fp panel title] at (0,2.02) {$n=31$};
  \node[fp active, minimum width=35mm, anchor=west] (a31) at (0,1.32)
    {$A=\bar Q_{15}$ and shift $s$\\outer square gadget};
  \node[fp seam, minimum width=8mm, anchor=west] at (3.5,1.32) {$-1$};
  \node[fp second, minimum width=27mm, anchor=west] (b31) at (4.3,1.32)
    {$B=Q_7$\\$xB^2$};
  \node[fp seam, minimum width=8mm, anchor=west] at (7.0,1.32) {$+1$};
  \node[fp pivot, minimum width=22mm, anchor=west] (h31) at (7.8,1.32)
    {$H_4$\\$H_4^2$};
  \node[fp seam, minimum width=8mm, anchor=west] at (10.0,1.32) {$-1$};
  \node[fp recursive, minimum width=28mm, anchor=west] (c31) at (10.8,1.32)
    {$C=Q_3$, then $r$\\low residual};
  \node[fp label, anchor=west] at (4.0,0.56)
    {reconstruct $H_2$ from $H_4$, then discharge all conditional decoders};
  \node[fp label, anchor=west, text=fpBlue!80!black] at (0,-0.03)
    {Subtract the red boundary contributions before the corresponding pivot; $b$ is already known at its seam.};
\end{tikzpicture}
}
  \caption{The three finite decoders in one visual template.  A coloured
  block denotes a monic polynomial recovered from the indicated square
  window; each red seam is a fixed boundary coefficient contributed by the
  next block.  Recovering a polynomial and decoding its internal parameters
  are deliberately separated: the lower annotations show when the required
  powers have finally become available and the corresponding conditional
  decoder may be discharged.}
  \label{fig:special-case-decoders}
\end{figure}

\begin{lemma}[Finite cost bases]\label{lem:special-cases-splittable}
    For each $n\in\{15,27,31\}$, over an $n$-admissible field the
    special-case construction for $n$ produces a splittable pair.
    Moreover, the direct descending decoders below give the required
    compatibility windows for these three finite constructions; no general
    composition result is needed.  The two components have a joint
    $(n-1)/2$-realization that records their monic quadratic and quartic
    byproducts.
    In particular, $P_{15}$, $P_{27}$, and $P_{31}$ are decodable.
\end{lemma}
\begin{proof}
    We give the finite decoders explicitly.  Write
    $\Phi_n=xT^{(1)}_n+T^{(2)}_n=P_n$ and $G_n=\rng n$.
    For any polynomial $W$ write $W_j=\coeff{W}{j}$, so
    $p_j=\coeff{\Phi_n}{j}$.  Every square root below is the descending
    monic-square recursion (so it uses only divisions by $2$), and all
    quantities described as \text{known} are coefficients of the displayed
    auxiliary powers or quantities recovered in an earlier row block.
    To make explicit why only the top half of a square is needed, if
    $W=x^e+\sum_{r<e}w_rx^r$, then
    \[
      \coeff{W^2}{2e-s}
        =2w_{e-s}+\sum_{r=1}^{s-1}w_{e-r}w_{e-(s-r)}
        \qquad(1\le s\le e).
    \]
    Thus the coefficients in degrees $2e,\ldots,e$ recover the monic
    polynomial $W$ successively.

    \paragraph{The degree-$15$ construction.}
    Put
    \[
      A=S^{(1)}_1,\quad H_2=x^2+bx+c,\quad
      U=H_2+\alpha _3=x^2+bx+d_3,\quad
      V=H_2+\alpha _2=x^2+bx+d_2,
    \]
    and write $H_4=x^4+h_3x^3+h_2x^2+h_1x+h_0$.  The identity
    \[
      P_{15}=(x+1)A^2+H_4^2-(x+1)U^2-V^2+(x+1)\alpha _1+\alpha _0
    \]
    is a relative square shell with $M=x+1$, $\deg A=7$, and error
    \[
      E_A=H_4^2-(x+1)U^2-V^2+(x+1)\alpha_1+\alpha_0.
    \]
    Here $\deg E_A\le8$ and $[x^8]E_A=1$.  Hence
    \Cref{lem:relative-square-shell} recovers $A$ from the top block, with
    the displayed $+1$ seam correction.  Set
    $R=\Phi_{15}-(x+1)A^2$.  Its next four
    coefficients give
    \[
      h_3=R_7/2,\qquad b=h_3/2,\qquad
      h_2=(R_6-h_3^2)/2,
    \]
    \[
      h_1=(R_5+1-2h_3h_2)/2,\qquad
      h_0=(R_4+2b+2-h_2^2-2h_3h_1)/2.
    \]
    The defining relation $H_4=H_2^2-(x+\alpha _5)^2+\alpha _4$
    then yields
    \[
      c=(h_2-b^2+1)/2,\quad
      \alpha _5=bc-h_1/2,\quad \alpha _4=h_0-c^2+\alpha _5^2.
    \]
    Finally put $R'=R-H_4^2$.  The four bottom equations are
    \[
      \begin{aligned}
      d_3&=-(R'_3+b^2+4b)/2,\\
      d_2&=-(R'_2+2b^2+2d_3+2bd_3)/2,\\
      \alpha _1&=R'_1+2bd_3+d_3^2+2bd_2,\\
      \alpha _0&=R'_0+d_3^2+d_2^2-\alpha _1.
      \end{aligned}
    \]
    Hence $\alpha _3=d_3-c$ and $\alpha _2=d_2-c$.  Once $H_2,H_4$
    are known, the unitriangular decoder for $Q_7$ (\Cref{lem:Q-unitriangular})
    recovers the parameters inside $A$; the side information is discharged
    by \Cref{lem:discharge-side-information} because $A,H_2,H_4$ have all
    been recovered from $P_{15}$.

    \paragraph{The degree-$27$ construction.}
    Write
    \[
      A=S^{(1)}_1,\quad C=S^{(2)}_1,\quad B=S^{(1)}_2,
      \quad H_2=x^2+ux+v,
    \]
    where $\deg A=13$, $\deg C=7$, and $\deg B=3$, and put
    $s=\alpha _1$ and $r=\alpha _0$.  The identity
    \[
      P_{27}=(x+1)A^2-(x+1)B^2+(x+1)s+C^2-H_2^2+r
    \]
    gives three consecutive relative square shells.  First set
    \[
      E_A=-(x+1)B^2+(x+1)s+C^2-H_2^2+r.
    \]
    Since $\deg E_A\le14$ and $[x^{14}]E_A=1$,
    \Cref{lem:relative-square-shell} with $M=x+1$ recovers $A$.  Set
    $R=\Phi_{27}-(x+1)A^2$.  Then
    \[
      R=C^2+E_C,\qquad
      E_C=-(x+1)B^2+(x+1)s-H_2^2+r,
    \]
    where $\deg E_C\le7$ and $[x^7]E_C=-1$.  The same lemma, now with
    $M=1$, recovers $C$.  Set $R'=R-C^2=E_C$.  After changing sign,
    \[
      -R'=(x+1)B^2+E_B,\qquad
      E_B=-(x+1)s+H_2^2-r,
    \]
    with $\deg E_B\le4$ and $[x^4]E_B=1$.  A third application with
    $M=x+1$ recovers $B$.  These are exactly the seams $+1,-1,+1$ in
    \Cref{fig:special-case-decoders}.  Finally, with
    $L=R'+(x+1)B^2=(x+1)s-H_2^2+r$, we have
    \[
      u=-L_3/2,\qquad v=-(L_2+u^2)/2,\qquad
      s=L_1+2uv,\qquad r=L_0+v^2-s.
    \]
    Decode $A=Q_{13}$ first using the recovered $H_2$; by
    \Cref{lem:Q4k+1-from-H2} this also reconstructs the quartic $H_4$ used
    by $C$.  Then decode $C=Q_7$ from $(H_2,H_4)$ and $B=Q_3$ from $H_2$.
    These three decoders recover their respective parameter blocks without
    assuming an internal power before it has been reconstructed; this is
    exactly the substitution in \Cref{lem:discharge-side-information}.

    \paragraph{The degree-$31$ construction.}
    Put
    \[
      A=S^{(1)}_1,\quad B=S^{(1)}_2,\quad C=S^{(1)}_3,
      \quad H=H_4,
    \]
    with degrees $15,7,3,4$, respectively, and set
    $s=\alpha _{15}$ and $r=\alpha _0$.  Since
    \[
      P_{31}=xA^2+(A+s)^2+E,
      \qquad E=-xB^2-H^2+xC+r,
    \]
    we have $\deg E\le15=\deg A$ and $[x^{15}]E=-1$.  Therefore
    \Cref{lem:square-gadget-boundary} recovers $A$ and $s$ directly, with
    the displayed fixed boundary correction.  Set
    \[
      R'=\Phi_{31}-xA^2-(A+s)^2.
    \]
    The remaining two monic blocks are again relative square shells:
    \[
      -R'=xB^2+E_B,\qquad E_B=H^2-xC-r.
    \]
    Here $\deg E_B\le8$ and $[x^8]E_B=1$, so
    \Cref{lem:relative-square-shell} with $M=x$ recovers $B$.  Put
    $R''=R'+xB^2$.  Then
    \[
      -R''=H^2+E_H,\qquad E_H=-xC-r,
    \]
    with $\deg E_H\le4$ and $[x^4]E_H=-1$.  The same lemma with $M=1$
    recovers $H$.  The remaining polynomial
    $R'''=R''+H^2=xC+r$ gives
    $r=R'''_0$ and $C=(R'''-r)/x$.  If
    $H=x^4+h_3x^3+h_2x^2+h_1x+h_0$, its parameter block is
    \[
      \alpha _7=h_3/2,\quad
      \alpha _6=(h_2-\alpha _7^2+1)/2,\quad
      \alpha _5=\alpha _7\alpha _6-h_1/2,\quad
      \alpha _4=h_0-\alpha _6^2+\alpha _5^2.
    \]
    The previously proved decoders for $Q_3,Q_7,\bar Q_{15}$ (the latter
    given $(H_2,H_4)$) recover all remaining parameters.  Formally,
    \Cref{lem:discharge-side-information} substitutes the recovered
    $H_2,H_4$ into these conditional decoders.

    It remains to record causality.  The preceding formulas have the following
    row shifts; $t$ denotes a coefficient degree inside the recovered
    polynomial.
    \[
    \begin{array}{c|c|c}
      n&\text{recovered quantity}&\text{row where its coefficient is read}\\ \hline
      15&[x^t]A,\ [x^t]H_4,\
          (\alpha_3,\alpha_2,\alpha_1,\alpha_0)
        &p_{t+8},\ p_{t+4},\ (p_3,p_2,p_1,p_0)\\
      27&[x^t]A,\ [x^t]C,\ [x^t]B,\ (H_2,s,r)
        &p_{t+14},\ p_{t+7},\ p_{t+4},\ (p_3,\ldots,p_0)\\
      31&[x^t]A,\ s,\ [x^t]B,\ [x^t]H_4,\ [x^t]C,\ r
        &p_{t+16},\ p_{15},\ p_{t+8},\ p_{t+4},\ p_{t+1},\ p_0
    \end{array}
    \]
    At a shared boundary the only contribution from the next block is the
    displayed monic coefficient $1$ or $-1$, which is subtracted before that
    block is entered (in the $31$ case this includes the row-$4$ correction
    for $xC$).

    Maintain the descending invariant that a quantity first solved in row
    $r$ is a polynomial in the observed rows $p_s$ with $s\ge r$ and in
    quantities solved above it.  Every displayed square or Cauchy recurrence
    has this form, so the invariant holds across all boundary rows.  The row
    shifts in the table then show that $[x^j]T^{(1)}$ uses only $p_i$ with
    $i\ge j+1$, whereas $[x^j]T^{(2)}$ uses only $p_i$ with $i\ge j$.
    These are exactly the cutoffs in \Cref{def:compatible-pair}.  Since the
    leading coefficient $p_n=1$ is fixed, all three pairs are compatible on
    $\rng n$, hence splittable, and their combined polynomials are decodable.

    Finally, the displayed algorithms are joint circuits for the two
    components.  Charging each shared intermediate only once gives
    \[
    \begin{array}{c|c|c}
      n&\text{products in the joint realization}&\text{total}\\ \hline
      15&H_2:1,\ H_4:1,\ Q_7:3,\ \text{two outer products}:2&7\\
      27&H_2:1,\ Q_{13}:6,\ Q_3:1,\ Q_7:3,\ \text{two outer products}:2&13\\
      31&H_2,H_4:2,\ \bar Q_{15}:7,\ Q_7:3,\ Q_3:1,\
          \text{two outer products}:2&15.
    \end{array}
    \]
    The $Q_{13}$ call in degree~$27$ records the required $H_4$; in the
    other two cases $H_2,H_4$ are displayed intermediates.  Hence all three
    realizations record both byproducts and have exactly $(n-1)/2$ products.
\end{proof}

